\documentclass[11pt,a4paper]{article}
\usepackage[utf8]{inputenc}
\usepackage[T1]{fontenc}
\usepackage[portuguese]{babel}
\usepackage{xcolor}
\usepackage{hyperref}
\usepackage{geometry}
\usepackage{tcolorbox}

\geometry{margin=2.5cm}

\definecolor{primary}{RGB}{41, 128, 185}
\definecolor{secondary}{RGB}{44, 62, 80}
\definecolor{codebg}{RGB}{240, 240, 240}

\hypersetup{colorlinks=true, linkcolor=primary, urlcolor=primary}

\title{\color{secondary}\Huge\textbf{Exercício 4.9: O projeto, etapa 25}}
\author{\color{primary}DevOps com Kubernetes -- 2026}
\date{}

\begin{document}
\maketitle

\section*{\color{primary}Instruções do Exercício}
\begin{tcolorbox}[colback=blue!5!white,colframe=blue!75!black,title=Exercício 4.9: O projeto, etapa 25]
Melhore as definições de \textit{deployment} do seu Projeto da seguinte forma com Kustomize e ArgoCD:

- Crie dois ambientes separados (\textit{production} e \textit{staging}) e os instancie em \textit{namespaces} diferentes no cluster.

- Cada commit aprovado na ramificação \texttt{main} deve acionar automaticamente o \textit{deploy} no ambiente de \textit{staging}.

- Cada commit que possuir uma \textbf{tag} de versão (\textit{release}) deve acionar o \textit{deploy} no ambiente de \textit{production}.

- No ambiente de \textit{staging}, o \textit{broadcaster} apenas deve registrar as mensagens em formato de logs, e não deverá enviá-las de verdade para os serviços de chat externos.

- No ambiente de \textit{staging}, não configure ou execute backups do banco de dados.
\end{tcolorbox}

\end{document}
